\documentclass[a4paper,12pt]{article}
\usepackage[utf8]{inputenc}
\usepackage[T1]{fontenc}
\usepackage[ngerman]{babel}
\usepackage{lmodern}
\usepackage{geometry}
\usepackage{hyperref}
\usepackage{booktabs}
\usepackage{enumitem}
\geometry{margin=25mm}
\hypersetup{
  colorlinks=true,
  linkcolor=black,
  urlcolor=blue,
  citecolor=black
}

\title{Zwischenbericht\\Campus Next-Gen Data-Hub}
\author{Lahres Timo \\ Bräutigam Rebecca \\ Horst Isabella}
\date{07.06.2026}

\begin{document}

\begin{titlepage}
  \centering
  {\LARGE \textbf{Zwischenbericht}}\\[1.5em]
  {\large \textbf{Campus Next-Gen Data-Hub}}\\[2em]
  {\large Evaluation von Airbyte als ETL- und Integrationswerkzeug}\\[0.5em]
  {\large Ablösung von Talend}\\[3em]
  {\large INF-M Modul Projekte SoSe '26}\\[0.5em]
  {\large Gruppe Airbyte}\\[2em]
  {\large Lahres Timo \\ Bräutigam Rebecca \\ Horst Isabella}\\[3em]
  {\large 07.06.2026}
\end{titlepage}

\tableofcontents
\newpage

\section{Projektziel \& Kontext}

Im Projekt evaluieren wir \href{https://airbyte.com/}{Airbyte} als modernes ETL- und Datenintegrationswerkzeug mit dem Ziel, das bisher eingesetzte \textbf{Talend} in der Hochschul-IT abzulösen.

Die Evaluationsumgebung läuft \textbf{lokal in Docker Desktop} und stellt einen realistischen Ausschnitt der Hochschul-Datenlandschaft dar, einschließlich anonymisierter Studierenden-, Gebäude-, Instituts- und Personaldaten.

Die Evaluation ist in \textbf{sechs Testszenarien} gegliedert:
\begin{itemize}[noitemsep]
  \item DB-Connector
  \item Facility-Management-Sync
  \item Bild-BLOBs
  \item Studenten/Personal-Mapping
  \item Incremental Sync / IdM
  \item Web-APIs
\end{itemize}

Siehe \href{https://github.com/Timbo3399/INFM_Airbyte/blob/main/docs/testszenarien.md}{testszenarien.md}.

\section{Erreichte Meilensteine (Stand 06.06.2026)}

\begin{tabular}{@{}p{6.5cm}p{4cm}p{4.5cm}@{}}
\toprule
\textbf{Meilenstein (lt. Betreuer-Mail)} & \textbf{Status} & \textbf{Beleg / Doku} \\
\midrule
Installation des Systems & ✅ DB-Stack + Airbyte (abctl) lauffähig & \href{https://github.com/Timbo3399/INFM_Airbyte/blob/main/docs/installation-guide.md}{installation-guide.md} \\
Zugang für Betreuer & ◑ dokumentiert, Einrichtung im Termin & \href{https://github.com/Timbo3399/INFM_Airbyte/blob/main/docs/zugang.md}{zugang.md} \\
Einfacher ETL-Prozess & ✅ \textbf{durchgeführt und verifiziert} (Postgres $\rightarrow$ Postgres) & \href{https://github.com/Timbo3399/INFM_Airbyte/blob/main/docs/etl-prozess.md}{etl-prozess.md} \\
Beginn der Dokumentation & ✅ README + 8 Dokumente unter \texttt{docs/} & dieses Repo \\
\bottomrule
\end{tabular}

\subsection{Installation}

Das Setup ist vollständig skriptbasiert und reproduzierbar:
\begin{itemize}[noitemsep]
  \item \texttt{scripts/install.ps1} (Windows) bzw. \texttt{scripts/install.sh} (Linux/macOS) startet den kompletten Datenbank-Stack (Source-PostgreSQL, Ziel-PostgreSQL, Ziel-MySQL, CSV-File-Server), wartet auf den \texttt{healthy}-Status und lädt die Testdaten.
  \item \texttt{scripts/setup-airbyte.ps1} / \texttt{scripts/setup-airbyte.sh} installiert Airbyte Community Edition über das offizielle CLI \texttt{abctl} in einem lokalen Kubernetes-Cluster (kind) innerhalb von Docker Desktop.
  \item \textbf{Plattformen:} Das Setup steht für \textbf{Windows, Linux und macOS} bereit (identische Logik, plattformspezifische Skripte).
\end{itemize}

Alle vier Datenbank-/Server-Container laufen verifiziert im Zustand \texttt{healthy}.

\subsection{Datenbasis (Source-PostgreSQL)}

Die anonymisierten Testdaten sind geladen:

\begin{tabular}{@{}p{5cm}r p{6cm}@{}}
\toprule
\textbf{Tabelle} & \textbf{Zeilen} & \textbf{Lademechanismus} \\
\midrule
\texttt{fm\_gebaeude} & 25 & \texttt{scripts/load\_fm\_gebaeude.py} \\
\texttt{fm\_inst} & 2.083 & \texttt{scripts/load\_fm\_inst.py} \\
\texttt{k\_plz} & 34.172 & \texttt{scripts/load\_k\_plz.py} \\
\texttt{fm\_rna} & 379 & \texttt{scripts/load\_json.py} \\
\texttt{hso\_personal} & 870 & \texttt{scripts/load\_json.py} \\
\texttt{hso\_students} & 0 & \textbf{nur via File-Connector (s. Kap. 4)} \\
\texttt{fm\_stamm} & 0 & keine Quelldatei (s. Kap. 5) \\
\bottomrule
\end{tabular}

\subsection{Angelegte Airbyte-Connectoren \& erster ETL-Lauf}

In Airbyte sind angelegt und per Verbindungstest erfolgreich getestet:
\begin{itemize}[noitemsep]
  \item \textbf{Sources (5):}
  \begin{itemize}[noitemsep]
    \item \texttt{HSO Source PostgreSQL} (Postgres, Update-Methode \textit{User Defined Cursor})
    \item \texttt{HSO CSV hso\_students} (\texttt{local}, \texttt{/local/*.csv})
    \item \texttt{HSO CSV k\_plz} (\texttt{local}, \texttt{/local/*.csv})
    \item \texttt{HSO CSV fm\_gebaeude} (\texttt{local}, \texttt{/local/*.csv})
    \item \texttt{HSO CSV fm\_inst} (\texttt{local}, \texttt{/local/*.csv})
  \end{itemize}
  \item \textbf{Destinations (2):}
  \begin{itemize}[noitemsep]
    \item \texttt{HSO Dest PostgreSQL} (Port 5434)
    \item \texttt{HSO Dest MySQL} (Port 3306, SSL aus, \texttt{allowPublicKeyRetrieval=true}, Raw-DB \texttt{destdb})
  \end{itemize}
\end{itemize}

Es wurden \textbf{drei Connections} (jeweils \textit{Full Refresh | Overwrite}) ausgeführt und das Ergebnis \textbf{unabhängig in der jeweiligen Ziel-DB} geprüft:

\begin{tabular}{@{}p{7cm}p{4cm}p{4cm}@{}}
\toprule
\textbf{Connection} & \textbf{Streams} & \textbf{Ergebnis (Ziel-DB)} \\
\midrule
\texttt{HSO Source PostgreSQL $\rightarrow$ HSO Dest PostgreSQL} & \texttt{fm\_gebaeude}, \texttt{k\_plz} & 25 / 34.172 ✅ \\
\texttt{HSO Source PostgreSQL $\rightarrow$ HSO Dest MySQL} & \texttt{fm\_gebaeude}, \texttt{k\_plz} & 25 / 34.172 ✅ \\
\texttt{HSO CSV hso\_students $\rightarrow$ HSO Dest PostgreSQL} & \texttt{hso\_students} (File) & \textbf{5.052 Zeilen} ✅ \\
\bottomrule
\end{tabular}

Bemerkenswert: Der \textbf{File-Connector lud die defekte \texttt{hso\_students.csv} vollständig (5.052 Zeilen)} in die DB — dieselbe Datei, an der ein direktes PostgreSQL-\texttt{COPY} scheiterte (0 Zeilen, Kap. 5.2). Pandas im File-Connector toleriert die Spalteninkonsistenzen. Damit sind \textbf{DB-Connector} (PG $\rightarrow$ PG, PG $\rightarrow$ MySQL) und \textbf{File-Connector} (CSV $\rightarrow$ DB) — also der Kern von Szenario 1 inkl. „PostgreSQL $\rightarrow$ MySQL dumpen“ — nachgewiesen.

Zusätzlich stellt der Dienst \textbf{PostgREST} (\texttt{hso\_postgrest}, Szenario 6a) eine REST-API auf die Ziel-DB bereit: \texttt{GET http://localhost:3000/k\_plz?limit=5} liefert die synchronisierten Daten als JSON.

\section{Architektur (Kurzüberblick)}

Detaillierte Beschreibung in \href{https://github.com/Timbo3399/INFM_Airbyte/blob/main/docs/architektur.md}{architektur.md}. Kern: alle Dienste laufen in Docker Desktop in einem gemeinsamen Netzwerk (\texttt{airbyte\_net}). Airbyte selbst läuft in einem kind-Cluster und erreicht die Datenbanken über \texttt{host.docker.internal}.

\begin{verbatim}
 Docker Desktop
 ┌──────────────────────────────────────────────────────────────┐
 │  source-postgres   ──┐                                        │
 │  (Testdaten)         │        Airbyte (kind-Cluster)          │
 │  localhost:5433      ├──────► UI  : localhost:8000            │
 │                      │        liest Source / schreibt Ziele   │
 │  file-server         │                                        │
 │  localhost:8888 ─────┘                │                       │
 │                          ┌────────────┴───────────┐           │
 │   dest-postgres  localhost:5434   dest-mysql  localhost:3306  │
 └──────────────────────────────────────────────────────────────┘
\end{verbatim}

\section{Besonderheit: Datenqualität der Quell-CSVs}

Die bereitgestellten CSV-Dateien ließen sich \textbf{nicht} per direktem PostgreSQL-\texttt{COPY} laden (Details in Kap. 5). Wir haben daher tolerante Python-Loader implementiert, die die Daten nach dem Containerstart bereinigt einspielen. Eine Datei — \texttt{hso\_students.csv} — ist strukturell so inkonsistent, dass sie aktuell nicht zuverlässig in die relationale Source-DB geladen werden kann; Studierendendaten werden stattdessen über den \textbf{Airbyte File-Connector} als Flatfile-Quelle eingebunden.

\section{Probleme \& Lösungen}

\subsection{Windows-/Umgebungsspezifische Hürden}

\begin{tabular}{@{}p{6cm}p{5cm}p{5cm}@{}}
\toprule
\textbf{Problem} & \textbf{Ursache} & \textbf{Lösung} \\
\midrule
\texttt{install.ps1} meldete „Python gefunden", JSON-Laden schlug aber fehl & \texttt{python} ist unter Windows oft nur der \textbf{Microsoft-Store-Platzhalter}; \texttt{Get-Command} findet ihn, er liefert aber keine echte Version & Echte Versionsprüfung (\texttt{py}/\texttt{python}/\texttt{python3}); zusätzlich \textbf{Docker-Fallback}, der die Daten ganz ohne Host-Python lädt \\
\texttt{setup-airbyte.ps1} fand kein abctl-Asset & Falsches Namensschema (\texttt{abctl-\_Windows\_amd64.zip}) — korrekt ist \texttt{abctl-<version>-windows-<arch>.zip}; zudem liegt \texttt{abctl.exe} in einem Unterordner & Asset per Muster ermitteln, aus Unterordner entpacken; Architektur-Erkennung PowerShell-7-fest \\
File-Server-Container dauerhaft \texttt{unhealthy}, Setup lief in Timeout & Healthcheck nutzte \texttt{localhost} → im Container zuerst IPv6 (\texttt{::1}), nginx lauscht aber nur auf IPv4 & Healthcheck auf \texttt{127.0.0.1} umgestellt \\
\bottomrule
\end{tabular}

\subsection{Datenqualität der Quell-CSVs (Hauptproblem)}

Der SQL-Init lud die drei Kern-Tabellen \textbf{gar nicht} — \texttt{COPY} brach unter \texttt{ON\_ERROR\_STOP} bereits an der ersten fehlerhaften Datei ab, wodurch alle folgenden leer blieben. Die konkreten Probleme:

\begin{itemize}[noitemsep]
  \item \textbf{\texttt{fm\_gebaeude.csv}}: unquotierte Kommas in Textfeldern (z. B. „Hörsäle,Bib.,RZ"), eingebettete Header-Zeilen.
  \item \textbf{\texttt{k\_plz.csv}}: \textbf{3.417 Header-Zeilen} mitten in den Daten (zusammengesetzte Einzel-Exporte), Zeilen mit zu wenig Feldern, ein Feld (\texttt{krskfz}) breiter als das Schema (Kreis-Namen).
  \item \textbf{\texttt{fm\_inst.csv}}: 86 Spalten (Tabelle nutzt 24), semikolon-getrennt, vereinzelt \textbf{NUL-Bytes}, doppelt kodierte UTF-8-Umlaute.
  \item \textbf{\texttt{hso\_students.csv}}: \textbf{strukturell defekt} — Datenzeilen haben mehr Spalten (~50–56) als der eigene Header (40), dazu unbalanciertes Quoting und Float-formatierte Integer.
\end{itemize}

\textbf{Lösung:} tolerante Python-Loader (\texttt{scripts/load\_*.py}), die Header filtern, fehlerhafte Zeilen reparieren bzw. überspringen, NUL-Bytes und Mojibake bereinigen und Typkonvertierungen best-effort vornehmen. Der SQL-\texttt{COPY}-Schritt wurde entfernt (\texttt{01\_load\_data.sql}).

\subsection{Airbyte / abctl-spezifische Hürden}

Airbyte Community läuft über \texttt{abctl} in einem \textbf{kind-Kubernetes-Cluster}. Daraus ergaben sich mehrere nicht-offensichtliche, in der offiziellen Doku nicht beschriebene Hürden, die wir analysiert und gelöst haben:

\begin{tabular}{@{}p{6cm}p{5cm}p{5cm}@{}}
\toprule
\textbf{Problem} & \textbf{Ursache} & \textbf{Lösung} \\
\midrule
\textbf{File-Connector} (\texttt{local}) findet \texttt{/local/*.csv} nicht & Connector-Pods (kind) sehen das Docker-Volume \texttt{oss\_local\_root} nicht & CSV-Verzeichnis beim Install via \texttt{abctl local install --volume "…:/local"} direkt in den Cluster mounten \\
\texttt{--volume} mit Windows-Pfad: \texttt{is not a valid volume spec} & abctl trennt den Volume-String stur an \texttt{:}, der Laufwerks-Doppelpunkt (\texttt{C:}) kollidiert & Pfad in MSYS-Form \texttt{/c/Users/...} angeben \\
\textbf{Alle} Connector-Tests/Syncs hängen \texttt{Pending} & Aktiviertes lokales Volume erwartet PVC \texttt{airbyte-local-pvc} in \textbf{jedem} Job-Pod; es existierte nicht (\texttt{persistentvolumeclaim not found}) & PV (hostPath \texttt{/local}) + PVC \texttt{airbyte-local-pvc} anlegen; \texttt{JOB\_KUBE\_LOCAL\_VOLUME\_ENABLED=true} + Neustart von launcher/worker \\
\texttt{abctl local credentials --email … --password …} schlägt fehl (\texttt{unable to determine organization email}, \texttt{invalid character '<'}) & Kombinierter Aufruf löst einen Org-Lookup aus, der HTML statt JSON liefert & E-Mail und Passwort in \textbf{zwei getrennten} Aufrufen setzen (erst \texttt{--email}, dann \texttt{--password}) \\
Connector-Auswahl heißt \textbf{„Postgres"} (nicht „PostgreSQL"); Default-Update-Methode ist \textbf{CDC} & — & In der UI „Postgres" wählen; Update-Methode auf \textit{User Defined Cursor} stellen (CDC bräuchte \texttt{wal\_level=logical}) \\
\bottomrule
\end{tabular}

Alle Lösungen sind in \texttt{scripts/setup-airbyte.ps1}/\texttt{.sh} automatisiert und in \href{https://github.com/Timbo3399/INFM_Airbyte/blob/main/docs/airbyte-setup.md}{airbyte-setup.md} / \href{https://github.com/Timbo3399/INFM_Airbyte/blob/main/docs/etl-prozess.md}{etl-prozess.md} dokumentiert.


\newpage

\section{Offene Punkte / Fragen an die Betreuer}

\begin{itemize}[noitemsep]
  \item \textbf{\texttt{hso\_students.csv} Soll-Struktur?} Die Datei ist nicht eindeutig ladbar (mehr Datenspalten als Header, defektes Quoting). Ist es gewollt, eine defekte Datei einzulesen, oder wird eine korrekt lesbare Datei bereitgestellt? Dies blockiert aktuell Szenario 4 (Account-Mapping in die Source-DB).
  \item \textbf{\texttt{fm\_stamm}} (Raumstammdaten): Es liegt keine Quelldatei vor. Wird diese Tabelle benötigt, und woher sollen die Daten kommen?
  \item \textbf{Zugang für Betreuer:} Airbyte Community Edition ist im Wesentlichen Single-User und läuft auf \texttt{localhost}. Wie soll Ihr Zugang erfolgen? Gemeinsamer Admin-Login während des Vor-Ort-Termins, oder ist ein Remote-Zugang (z. B. Tunnel/VM) gewünscht? (Vorschlag in \href{https://github.com/Timbo3399/INFM_Airbyte/blob/main/docs/zugang.md}{zugang.md}.)
  \item \textbf{Sync-Strategie:} Wir nutzen den Cursor-Modus (\texttt{updatedat}) statt CDC, da CDC zusätzliche PostgreSQL-Konfiguration (logical WAL, Replication Slot) erfordert. Ist CDC für die Evaluation gewünscht?
  \item \textbf{Scope:} Welche der sechs Szenarien haben für die Bewertung Priorität?
\end{itemize}

\section{Anforderungs- und Szenarien-Status}

Eine vollständige Gegenüberstellung aller Kickoff-Anforderungen und der sechs Szenarien mit Bewertung der Airbyte-Eignung und unserem Umsetzungsstand steht in \href{https://github.com/Timbo3399/INFM_Airbyte/blob/main/docs/anforderungen.md}{anforderungen.md}. Kernbefunde:

\begin{itemize}[noitemsep]
  \item \textbf{Erfüllt:} Open Source/Community, PostgreSQL- \& MySQL-Anbindung, CSV/JSON/Excel-Dateien, Logging/Monitoring, einfacher ETL-Lauf (verifiziert).
  \item \textbf{Einschränkungen ggü. Talend:} \textbf{kein} OSS-Connector für \textbf{Informix}; \textbf{kein XML} nativ; \textbf{keine freie Code-Snippet-Ausführung} (Mapping nur via dbt-SQL/Custom-Connector).
\end{itemize}

Diese Punkte sind die zentralen Evaluationsbefunde.

\section{Nächste Schritte}

\begin{itemize}[noitemsep]
  \item Szenario 1 abrunden: Sync auch nach MySQL + ein File $\rightarrow$ DB-Sync (Nachweis File-Connector-Last).
  \item Weitere Szenarien: FM-Denormalisierung (dbt/View), BLOB-Bilder, Incremental Sync (IdM), Web-APIs (PostgREST/SOAP).
  \item Klärung der offenen Fragen aus Kap. 6.
\end{itemize}

\section{Anhang: Repository-Struktur}

Siehe \href{https://github.com/Timbo3399/INFM_Airbyte/blob/main/README.md}{README.md}. Setup in unter 20 Minuten via \texttt{scripts/install.ps1} + \texttt{scripts/setup-airbyte.ps1}.

\end{document}
